\documentclass[11pt, a4paper]{article}

% ============================================================
% Packages
% ============================================================
\usepackage[utf8]{inputenc}
\usepackage[T1]{fontenc}
\usepackage{lmodern}
\usepackage[margin=1in]{geometry}
\usepackage{booktabs}
\usepackage{longtable}
\usepackage{float}
\usepackage{graphicx}
\usepackage{xcolor}
\usepackage{hyperref}
\usepackage{fancyhdr}
\usepackage{titlesec}
\usepackage{array}
\usepackage{caption}
\usepackage{amsmath}           % 解决 \text 未定义的问题

% ============================================================
% Page style
% ============================================================
\pagestyle{fancy}
\fancyhf{}
\fancyhead[L]{\small\textit{Credit Portfolio Risk Report}}
\fancyhead[R]{\small\textit{2026-04-03 11:06:56}}
\fancyfoot[C]{\thepage}
\renewcommand{\headrulewidth}{0.4pt}
\renewcommand{\footrulewidth}{0.4pt}

% ============================================================
% Title formatting
% ============================================================
\titleformat{\section}{\Large\bfseries\color{blue!60!black}}{\thesection}{1em}{}
\titleformat{\subsection}{\large\bfseries\color{blue!40!black}}{\thesubsection}{1em}{}

% ============================================================
% Hyperref
% ============================================================
\hypersetup{
    colorlinks=true,
    linkcolor=blue!70!black,
    citecolor=blue!70!black,
    urlcolor=blue!70!black,
}

% ============================================================
\begin{document}

% --- Title Page ------------------------------------------------
\begin{titlepage}
\centering
\vspace*{2cm}

{{\Huge\bfseries\color{blue!60!black} Credit Portfolio Risk Report\par}}

\vspace{0.5cm}
{{\Large Lending Club Loan Portfolio Analytics\par}}

\vspace{1.5cm}

\begin{tabular}{r l}
\textbf{Portfolio:}   & 2,260,668 Loans \\
\textbf{Exposure:}    & \$34,004,208,600 \\
\textbf{Grades:}      & A -- G \\
\textbf{Data Source:} & Lending Club + US Census ACS S2503 \\
\textbf{Generated:}   & 2026-04-03 11:06:56 \\
\end{tabular}

\vspace{1.5cm}

\begin{table}[H]
\centering
\caption{Executive Summary}
\begin{tabular}{l r}
\toprule
\textbf{Metric} & \textbf{Value} \\   % 表头加粗
\midrule
Total Loans & 2,260,668 \\
Total Funded Amount & \$34,004,208,600 \\
Average Interest Rate & 17.91\% \\
Mature Loans & 1,348,099 \\
Total Defaults & 269,360 \\
Overall Default Rate & 19.98\% \\
Expected Loss & \$3,107,705,485 \\
EL as \% of Exposure & 9.14\% \\
Average LGD & 46.22\% \\
HHI Index & 6,112 / 10,000 \\
\bottomrule
\end{tabular}
\label{tab:summary}
\end{table}

\vfill
{\small\textit{This report was auto-generated by the Credit Portfolio Risk Analytics Engine.}}
\end{titlepage}

% --- Table of Contents -----------------------------------------
\newpage
\tableofcontents
\newpage

% ============================================================
\section{Portfolio Overview}
\label{sec:portfolio}

The portfolio consists of \textbf{2,260,668} loans with a total funded
amount of \textbf{\$34,004,208,600}.  The average interest rate across all
grades is \textbf{17.91\%}.  The table below breaks down key
metrics by credit grade, from A (lowest risk) to G (highest risk).  Higher grades command
larger average loan amounts but carry substantially higher default rates, as shown in
subsequent sections.

\begin{table}[H]
\centering
\caption{Portfolio Overview by Grade}
\begin{tabular}{c r r r r r}
\toprule
Grade & Count & Total Funded (\$) & Avg Amount (\$) & Avg Rate (\%) & Avg FICO \\
\midrule
A & 433,027 & 6,321,329,325 & 14,598 & 7.08 & 731 \\
B & 663,557 & 9,401,557,325 & 14,168 & 10.68 & 702 \\
C & 650,053 & 9,773,557,750 & 15,035 & 14.14 & 691 \\
D & 324,424 & 5,096,087,575 & 15,708 & 18.14 & 686 \\
E & 135,639 & 2,365,386,850 & 17,439 & 21.83 & 684 \\
F & 41,800 & 798,480,275 & 19,102 & 25.45 & 682 \\
G & 12,168 & 247,809,500 & 20,366 & 28.07 & 681 \\
\textbf{TOTAL} & \textbf{2,260,668} & \textbf{34,004,208,600} & \textbf{15,042} & \textbf{17.91} & - \\
\bottomrule
\end{tabular}
\label{tab:portfolio}
\end{table}

% ============================================================
\section{Default Analysis}
\label{sec:defaults}

Default rates are computed on \emph{mature loans only} --- those that have reached a terminal
status (Fully Paid or Charged Off / Default).  Of the \textbf{1,348,099}
mature loans, \textbf{269,360} have defaulted, yielding an overall
default rate of \textbf{19.98\%}.  Default rates increase monotonically
from Grade A to Grade G, reflecting the underlying credit risk differentiation embedded in
Lending Club's grading methodology.

\begin{table}[H]
\centering
\caption{Default Rate by Grade (Mature Loans Only)}
\begin{tabular}{c r r r}
\toprule
Grade & Mature Count & Defaults & Default Rate (\%) \\
\midrule
A & 235,193 & 14,214 & 6.04 \\
B & 393,102 & 52,661 & 13.40 \\
C & 382,323 & 85,805 & 22.44 \\
D & 201,657 & 61,264 & 30.38 \\
E & 94,192 & 36,199 & 38.43 \\
F & 32,306 & 14,585 & 45.15 \\
G & 9,326 & 4,632 & 49.67 \\
\textbf{TOTAL} & \textbf{1,348,099} & \textbf{269,360} & \textbf{19.98} \\
\bottomrule
\end{tabular}
\label{tab:defaults}
\end{table}

% ============================================================
\section{Loss Given Default}
\label{sec:lgd}

LGD measures the fraction of exposure that is \emph{lost} when a loan defaults.  It is
calculated as:

\begin{equation}
    \text{LGD} = 1 - \frac{\text{total\_pymnt}}{\text{funded\_amnt}}
\end{equation}

The average LGD across all grades is \textbf{46.2\%}.  Lower grades
tend to exhibit higher LGD, driven by higher interest rates accumulating in the balance
before the borrower ultimately defaults.

\begin{table}[H]
\centering
\caption{Loss Given Default Statistics by Grade}
\begin{tabular}{c r r r r r}
\toprule
Grade & Count & Avg LGD (\%) & P25 (\%) & Median (\%) & P75 (\%) \\
\midrule
A & 14,214 & 44.9\% & 23.2\% & 46.0\% & 66.2\% \\
B & 52,661 & 44.6\% & 24.6\% & 46.7\% & 65.6\% \\
C & 85,805 & 46.2\% & 27.5\% & 49.7\% & 67.4\% \\
D & 61,264 & 46.3\% & 27.0\% & 50.4\% & 68.6\% \\
E & 36,199 & 45.8\% & 26.6\% & 50.7\% & 68.4\% \\
F & 14,585 & 46.0\% & 27.7\% & 52.0\% & 69.5\% \\
G & 4,632 & 49.7\% & 33.3\% & 56.5\% & 72.1\% \\
\bottomrule
\end{tabular}
\label{tab:lgd}
\end{table}

% ============================================================
\section{Expected Loss}
\label{sec:el}

Expected Loss (EL) is the risk-adjusted anticipated loss on the portfolio, computed per
grade segment as:

\begin{equation}
    \text{EL} = \text{PD} \times \text{LGD} \times \text{EAD}
\end{equation}

The total portfolio expected loss is \textbf{\$3,107,705,485}, representing
\textbf{9.14\%} of total exposure.  This is a key input for loan
pricing, capital allocation, and IFRS\,9 provisioning.

\begin{table}[H]
\centering
\caption{Expected Loss Breakdown by Grade}
\begin{tabular}{c r r r r r r}
\toprule
Grade & Exposure (\$) & PD & LGD & EL / Loan (\$) & Total EL (\$) & EL (\%) \\
\midrule
A & 6,321,329,325 & 6.0\% & 44.9\% & 395 & 171,241,019 & 0.50\% \\
B & 9,401,557,325 & 13.4\% & 44.6\% & 848 & 562,378,595 & 1.65\% \\
C & 9,773,557,750 & 22.4\% & 46.2\% & 1,559 & 1,013,690,735 & 2.98\% \\
D & 5,096,087,575 & 30.4\% & 46.3\% & 2,209 & 716,812,621 & 2.11\% \\
E & 2,365,386,850 & 38.4\% & 45.8\% & 3,071 & 416,512,124 & 1.22\% \\
F & 798,480,275 & 45.1\% & 46.0\% & 3,969 & 165,908,471 & 0.49\% \\
G & 247,809,500 & 49.7\% & 49.7\% & 5,026 & 61,161,920 & 0.18\% \\
\textbf{TOTAL} & \textbf{34,004,208,600} & - & - & - & \textbf{3,107,705,485} & \textbf{9.14\%} \\
\bottomrule
\end{tabular}
\label{tab:el}
\end{table}

% ============================================================
\section{Portfolio Concentration}
\label{sec:concentration}

The Herfindahl--Hirschman Index (HHI) measures portfolio concentration.  An HHI below
1,500 indicates a well-diversified portfolio; 1,500--2,500 is moderately concentrated;
above 2,500 is highly concentrated.  The total HHI for this portfolio is
\textbf{6,112}, indicating a \textbf{Highly Concentrated} profile.

\begin{table}[H]
\centering
\caption{Portfolio Concentration (HHI Index) -- Highly Concentrated}
\begin{tabular}{l l r r}
\toprule
Segment & Type & Share (\%) & HHI Contribution \\
\midrule
debt\_consolidation & purpose & 56.53 & 3,195 \\
B & grade & 29.35 & 862 \\
C & grade & 28.75 & 827 \\
credit\_card & purpose & 22.87 & 523 \\
A & grade & 19.15 & 367 \\
D & grade & 14.35 & 206 \\
home\_improvement & purpose & 6.66 & 44 \\
other & purpose & 6.17 & 38 \\
E & grade & 6.00 & 36 \\
major\_purchase & purpose & 2.23 & 5 \\
F & grade & 1.85 & 3 \\
medical & purpose & 1.22 & 1 \\
\textbf{TOTAL} &  &  & \textbf{6,112} \\
\bottomrule
\end{tabular}
\label{tab:concentration}
\end{table}

% ============================================================
\section{Interest Rate Analytics}
\label{sec:rates}

Interest rates vary significantly across credit grades, reflecting the risk-based pricing
model.  Grade A loans average 7.08\% while Grade G loans can reach up to 30.99\%.  This spread compensates investors for the higher probability of default in
lower-grade segments.

\begin{table}[H]
\centering
\caption{Interest Rate Distribution by Grade}
\begin{tabular}{c r r r r}
\toprule
Grade & Count & Min (\%) & Avg (\%) & Max (\%) \\
\midrule
A & 433,027 & 5.31 & 7.08 & 9.63 \\
B & 663,557 & 6.00 & 10.68 & 14.09 \\
C & 650,053 & 6.00 & 14.14 & 17.27 \\
D & 324,424 & 6.00 & 18.14 & 22.35 \\
E & 135,639 & 6.00 & 21.83 & 27.27 \\
F & 41,800 & 6.00 & 25.45 & 30.75 \\
G & 12,168 & 6.00 & 28.07 & 30.99 \\
\bottomrule
\end{tabular}
\label{tab:rates}
\end{table}

% ============================================================
\section{Model Performance}
\label{sec:models}

The following machine-learning models have been trained on the portfolio data:

\begin{itemize}
    \item \textbf{PD Model} --- HistGradientBoostingClassifier on 1,346,111 mature loans
          (AUC-ROC = 0.7235).
    \item \textbf{LGD Model} --- GradientBoostingRegressor with Huber loss on 269,360
          defaulted loans ($R^{2}$ = 0.1513).
    \item \textbf{Vasicek ASRF} --- Monte Carlo simulation (100,000 scenarios) with
          grade-grouped binomial approximation.  Expected Loss = \$3.1B,
          VaR @ 99.9\% = \$10.1B.
\end{itemize}

\subsection{PD Model}

The PD model achieves an AUC-ROC of \textbf{0.7235}, indicating good discriminative power
between defaulting and non-defaulting loans.  The model was trained on 1,076,888 loans
and validated on 269,223 holdout loans.  Actual vs.\ predicted default rates by grade show
strong calibration, with per-grade errors below 0.15 percentage points.

\subsection{LGD Model}

The LGD model predicts recovery rates with an $R^{2}$ of \textbf{0.1513} and MAE of
\textbf{0.2097}.  The top predictive feature is interest rate (importance = 0.47), followed
by sub-grade (0.12) and revolving utilisation (0.06).

\subsection{Vasicek ASRF Model}

The Vasicek model estimates the portfolio loss distribution under macroeconomic stress.
Key results:

\begin{itemize}
    \item Expected Loss: \$3,114,079,959 (9.16\% of exposure)
    \item VaR @ 99.9\%: \$10,080,224,115 (29.64\% of exposure)
    \item Unexpected Loss (UL): \$6,966,144,156 (20.49\%)
    \item Economic Capital Requirement: \$10,080,224,115
\end{itemize}

Asset correlations follow Basel II corporate exposure parameters:
$\rho_{A-B} = 0.15$, $\rho_{C-D} = 0.20$, $\rho_{E-G} = 0.25$.

% ============================================================
\end{document}